\documentclass[10pt,a4paper]{amsart}
\usepackage[margin=19mm]{geometry}
\usepackage{amsmath,amssymb,mathtools}
\usepackage[scr=rsfs,cal=euler]{mathalfa}
\usepackage{xfrac}
\usepackage[unicode,hidelinks,bookmarks=false]{hyperref}
\newcommand{\ie}{i.e.\ }
\newcommand{\cf}{cf.\ }
\newcommand{\resp}{respectively\ }
\newcommand{\Subsection}{\subsection}
\providecommand{\nobreakdash}{\nobreak\hbox{-}}
\setlength{\emergencystretch}{2em}
\multlinegap0pt
\usepackage{amssymb}
\usepackage{xcolor}
\newtheorem{lem}{Lemma}
\newcommand{\HH}{\mathrm{HH}}
\renewcommand{\L}{\mathbb{L}}
\newcommand{\Poly}{\mathrm{Poly}}
\title{Syntomic cohomology and real topological cyclic homology}
\author{Doosung Park}
\date{Source: arXiv:2311.06593v3 (CC BY 4.0)}
\begin{document}
\maketitle

\noindent\emph{Source and licence.} Syntomic cohomology and real topological cyclic homology. Version arXiv:2311.06593v3 (CC BY 4.0), distributed by arXiv under the Creative Commons Attribution 4.0 International licence, http://creativecommons.org/licenses/by/4.0/, as recorded in the arXiv licence metadata for this exact version and shown on the abstract page https://arxiv.org/abs/2311.06593v3. Full licence notice: \texttt{licenses/arxiv-2311.06593-CC-BY-4.0.txt}.\par\smallskip

\noindent\textit{Editorial scope.} Counted unit: the article's lem conv.1 (source0.tex lines 2260--2269). The article's complete original proof of it is delivered with it (source0.tex lines 2270--2342, one \texttt{proof} environment of 73 lines). No mathematics is written, repaired or re-derived: the statement and the proof are the article's own lines, in the article's own notation, and no retained line is substituted or re-flowed.  No further source-local context is needed: every cross-reference of every retained line resolves inside the two delivered spans.  Nothing was omitted from any retained span, so the package carries no omission record.  No retained line contains a citation, so the wrapper carries no bibliography and no bibliography entry.  No retained line contains a figure environment, an image-inclusion command or drawing code, so no image is delivered and the package declares no asset dependency.  Editorial formatting only, in addition: an AMS \texttt{amsart} preamble that restates the article's own declarations and macros used by the retained lines, namely the environments \texttt{lem} on separate counters, together with the standard packages those lines need.  The retained environments print in this document as Lemma 1 in order of appearance, whereas the article prints the same items with the numbering of its own full document.  The complete native source of the article as distributed in the arXiv e-print for this exact version is bundled unchanged as \texttt{sources/149876/source0.tex}.
\begin{lem}
\label{conv.1}
Let $R$ be a commutative ring of characteristic $2$.
For $A\in \Poly_R$,
there is a natural equivalence
\[
\HH(A/A^{(1)})\simeq
(\Gamma_A(\Omega_{A/R}^1),0).
\]
\end{lem}

\begin{proof}
Assume $A=R[x_1,\ldots,x_n]$.
Since $A^{(1)}\to A$ is flat,
we have
\[
B:=A\otimes_{A^{(1)}}A\simeq A\otimes_{A^{(1)}}^\L A.
\]
We can write the divided power algebra $\Gamma_B(\Omega_{A/R}^1\otimes_A B)$ as the free $B$-module
\[
\Gamma_B(\Omega_{A/R}^1\otimes_A B)
\cong
\bigoplus_{i_1,\ldots,i_n\geq 0} B \gamma_{i_1}(x_1)\cdots \gamma_{i_n}(x_n).
\]
The multiplication on  $\Gamma_B(\Omega_{A/R}^1\otimes_A B)$ is given by the formula
\[
\gamma_i(x_r)\gamma_j(x_r)
:=
\binom{i+j}{i}\gamma_{i+j}(x_r)
\]
for $1\leq r\leq n$.
Consider the map $d\colon \Gamma_B(\Omega_{A/R}^1\otimes_A B)\to \Gamma_B(\Omega_{A/R}^1\otimes_A B)$ given by the formula
\begin{align*}
&d(y\gamma_{i_1}(x_1)\cdots \gamma_{i_n}(x_n)) 
\\
:= &
\epsilon_1 y\gamma_{i_1-1}(x_1)\cdots \gamma_{i_n}(x_n)
+
\cdots
+
\epsilon_n y\gamma_{i_1}(x_1)\cdots \gamma_{i_n-1}(x_n)
\end{align*}
for $y\in B$,
where $\epsilon_r:=x_r\otimes 1-1\otimes x_r$ for $1\leq r\leq n$,
and $\gamma_{-1}=0$.
Observe that $d$ is a derivation.
We can form the commutative differential graded algebra $(\Gamma_B(\Omega_{A/R}^1\otimes_A B),d)$.

We claim that the truncation map
\begin{equation}
\label{conv.1.1}
(\Gamma_B(\Omega_{A/R}^1\otimes_A B),d)
\to
A
\end{equation}
is a quasi-isomorphism of commutative differential graded algebras.
If $n=1$,
then we can check this using $\epsilon_1^2=0$.
For general $n$,
we set $A_i:=R[x_i]$ and $B_i:=A_i\otimes_{A_i^{(1)}}A_i$.
Then there is an isomorphism of commutative differential graded algebras
\[
(\Gamma_B(\Omega_{A/R}^1\otimes_A B),d)
\cong
(\Gamma_{B_1}(\Omega_{A_1/R}^1\otimes_{A_1} B_1),d)
\otimes_R
\cdots
\otimes_R
(\Gamma_{B_n}(\Omega_{A_n/R}^1\otimes_{A_n} B_n),d).
\]
Using the case of $n=1$,
we see that the right-hand side is quasi-isomorphic to $A_1\otimes_R \cdots \otimes_R A_n$,
which is isomorphic to $A$.

Hence we have a quasi-isomorphism
\[
(\Gamma_B(\Omega_{A/R}^1\otimes_A B),d)\otimes_B A
\simeq
A\otimes_B^\L A.
\]
The left hand-side is isomorphic to $(\Gamma_A(\Omega_{A/R}^1),0)$ since the image of $\epsilon_r$ in $A$ is $0$ for $1\leq r\leq n$.
Hence we obtain the desired equivalence,
which is natural in $A$ since the truncation map \eqref{conv.1.1} is natural in $A$.
\end{proof}

\end{document}
